\subsection{A critical covering condition with a subpower margin}
\label{endpoint:covering-section}

The strict inequality in Proposition~\ref{prop:coherentcovering} can be
replaced by a quantitative covering condition at its critical exponent.
This does not replace the dimension inequality by a non-strict one for
arbitrary sets. The extra information concerns the number of occupied
squares at each scale and is invisible to the value of packing dimension.

Let \(1<s<2\), let \(\mu\) be an \(s\)-Frostman probability on a
compact planar set \(K\), and let \(N_n\) count the occupied dyadic
squares of side \(2^{-n}\). Changing the convention at square boundaries
changes the relevant covering estimates by at most a fixed factor.
First take a compactly supported \(t\)-Frostman pin probability \(\nu\),
\(t>1\), with support positively separated from \(K\). Use the same
angular exponent \(q>1\), centers, and positive affine densities \(f_n\)
as above, and put \(\eta_q=(q-1)/(2q-1)\).

\begin{proposition}[Critical covering summability]
\label{endpoint:criterion}
If
\begin{equation}\label{endpoint:series}
 \sum_{n\ge1}
 \left(n\,2^{-n(2s-1)}N_{n+1}\right)^{\eta_q}<\infty,
\end{equation}
then the raw joint distance law of \(\mu\times\nu\) is absolutely
continuous relative to \(d\nu(y)dt\).
\end{proposition}

\begin{proof}
For all sufficiently large \(n\), set
\[
 r=2^{-n},\qquad a_n=s-\frac1n,\qquad
 \gamma_n=\frac{a_n-1}{2}.
\]
Then
\[
 \frac{s-1}{4}\le\gamma_n<\frac{s-1}{2}<\frac12.
\]
All constants in Lemma~\ref{lem:coherentangular} and
Proposition~\ref{prop:coherentcomparison} can therefore be chosen
uniformly in \(n\). To justify this use explicitly, the Riesz identity,
angular interpolation, the angular Slobodeckij integral, and the scalar
translation bound have constants bounded when their exponent ranges over
this compact subinterval of \((0,1/2)\). The energy orders
\(1+2\gamma_n\) stay in a compact subinterval of \((1,2)\).
The geometric constants and the radial-projection exponent \(q\)
are fixed. The only constant not uniform in the ensuing Frostman-energy
calculation is the explicit factor \((s-a_n)^{-1}=n\).

The centers and affine densities were defined without reference to
\(\gamma\). Consequently we compare the same consecutive functions
\(f_{n+1},f_n\) at each step, using \(\gamma_n\) only in their estimate.
All exceptional angular sets remain null after taking the countable union
over nodes and over these exponents.

For a child square \(Q\) of diameter at most \(C r\), the Frostman
bound and layer cake give, uniformly for \(x\in Q\),
\begin{align*}
 \int_Q|x-x'|^{-a_n}\,d\mu(x')
 &\le (Cr)^{-a_n}\mu(Q)
       +a_n\int_0^{Cr}\mu(B(x,\rho))\rho^{-a_n-1}\,d\rho\\
 &\le C\left(1+\frac1{s-a_n}\right)r^{s-a_n}
 \le C n r^{1/n}.
\end{align*}
Integrating in \(x\) and normalizing twice shows
\[
 p_Q I_{a_n}(\mu_Q)\le C n r^{1/n}.
\]
The quantity in the square root of the comparison, apart from its
angular threshold, is therefore bounded by
\begin{align*}
 Z_n^*
 &:=r^{1+4\gamma_n}\sum_{Q\text{ at level }n+1}
                      p_Q I_{a_n}(\mu_Q)\\
 &\le C n N_{n+1} r^{2s-1-1/n}
 = C' n\,2^{-n(2s-1)}N_{n+1}.
\end{align*}
Here \(r^{-1/n}=2\), exactly.

The comparison, uniformly in \(n\), is
\[
 \|f_{n+1}-f_n\|_1
 \le C\bigl[(A Z_n^*)^{1/2}+B A^{1-q}\bigr]
 \qquad(A\ge1).
\]
The hypothesis makes \(Z_n^*\to0\). Optimize \(A\) exactly as in
Theorem~\ref{thm:coherentcriterion}, with \(B\) enlarged to at least one.
It gives
\[
 \|f_{n+1}-f_n\|_1
 \le C B^{1/(2q-1)}(Z_n^*)^{\eta_q}.
\]
The series of differences converges. Its nonnegative limit has mass one,
and the uniform affine approximation identifies it with the actual joint
distance law. This proves the assertion.
\end{proof}

For example, a bound
\(N_n\le C2^{n(2s-1)}n^{-b}\) is sufficient for this fixed separated
pair whenever \(b>1+1/\eta_q\). This is only a sufficient logarithmic
power; no optimality of that exponent is asserted.

\begin{theorem}[A positive-distance class at the critical dimension]
\label{endpoint:subpower}
Let \(1<s<2\), and suppose an \(s\)-Frostman probability \(\mu\)
is supported on a compact planar set \(K\) satisfying, for all
sufficiently large integers \(n\),
\begin{equation}\label{endpoint:subpower-count}
 N_n\le C\,2^{n(2s-1)}\exp(-c\sqrt n),\qquad C,c>0.
\end{equation}
For every compactly supported \(t\)-Frostman probability \(\nu\),
\(t>1\), its raw joint distance law with source \(\mu\) is absolutely
continuous. The supports need not be separated. In particular, taking
\(\nu=\mu\),
\[
 |\Delta_y(K)|>0\qquad\text{for \(\mu\)-almost every }y\in K.
\]
\end{theorem}

\begin{proof}
For separated supports, \eqref{endpoint:subpower-count} bounds the
summands in \eqref{endpoint:series} by
\(C n^{\eta_q}\exp(-c'\eta_q\sqrt n)\), a summable sequence for
every \(\eta_q>0\). Thus no quantitative lower bound for the radial
exponent \(q-1\) is needed.

For general compact supports cover the complement of the source--pin
diagonal by countably many products of positively separated closed bounded
balls. Normalize the restrictions on every product of positive measure.
The source restriction remains \(s\)-Frostman and its support has at most
a fixed multiple of the original covering count; the pin restriction
remains \(t\)-Frostman. Hence the separated result applies, even though
its constants and angular exponent may depend on the product. The diagonal
has zero pair measure since \(\mu\) is nonatomic. A reference-null joint
set has a null preimage on every product, and consequently under the full
pair probability. This proves joint absolute continuity without separation.
Disintegration and the positive total pinned mass give the final assertion.
\end{proof}

The numerical decay \(\exp(-c\sqrt n)\) is convenient rather than
canonical. The summability condition \eqref{endpoint:series} is the actual
criterion used. In particular the theorem concerns additional covering
information at \(\overline{\dim}_{\rm B}K\le2s-1\); the dimension
equality alone does not imply that information.
